\documentclass[12pt]{article}
\usepackage[T1]{fontenc}
\usepackage[utf8]{inputenc}
\usepackage{lmodern}
\usepackage{textcomp}
\usepackage{enumitem}
\usepackage{geometry}
\geometry{margin=1in}
\begin{document}
\begin{center}
Answer Key - History - K-12 Level Assessment 21 \
\small (Answers are ordered exactly as in the question paper.)
\end{center}

\section*{Multiple Choice}
\begin{enumerate}[label=\arabic*.]
\item \textbf{Answer:} D
\\ \textit{Options:} A) Surplus food is stored in large, collective facilities.; B) Crop irrigation requires the construction of dams or canals.; C) Sedentary communities need to defend themselves.; D) All of the above.
\\ \textit{Note:} The correct answer "All of the above" is appropriate because societies become more complex through a combination of factors such as increased population, technological advancements, and the development of social hierarchies, which collectively enhance interactions and interdependencies within the community.
\item \textbf{Answer:} D
\\ \textit{Options:} A) new and improved stone tool technologies.; B) a broadening of the subsistence base.; C) larger sites with increased populations.; D) all of the above.
\\ \textit{Note:} The Upper Paleolithic is characterized by significant advancements in human culture, including the development of art, sophisticated tools, and complex social structures, which collectively encompass all the options presented in the question.
\item \textbf{Answer:} C
\\ \textit{Options:} A) argon/argon and potassium argon dating; B) archaeomagnetic and radiocarbon dating; C) luminescence dating and optically stimulated luminescence; D) X-ray fluorescence and paleomagnetic dating
\\ \textit{Note:} Luminescence dating and optically stimulated luminescence are correct because they measure the energy released from minerals, such as quartz and feldspar, that have been exposed to heat and light, allowing researchers to determine when the material was last heated or exposed to sunlight.
\item \textbf{Answer:} C
\\ \textit{Options:} A) natural selection; B) catastrophism; C) uniformitarianism; D) unilineal evolution
\\ \textit{Note:} Uniformitarianism refers to the geological principle that the processes shaping the Earth today, such as erosion and sedimentation, have been consistent over time, encapsulated in the phrase "the slow agency of existing causes," which emphasizes the gradual and observable nature of these changes.
\item \textbf{Answer:} B
\\ \textit{Options:} A) the site of An-yang near the Huang ho River.; B) chiefdoms and states in numerous regions throughout China.; C) the beginning of the Qin Dynasty.; D) external influences originating from the Indus Valley.
\\ \textit{Note:} The correct answer highlights that Chinese civilization emerged from a variety of chiefdoms and states, indicating the early complexity and regional diversity in social and political organization that laid the foundation for the rich history of China.
\end{enumerate}

\section*{True / False}
\begin{enumerate}[label=\arabic*.]
\setcounter{enumi}{5}
\item \textbf{Answer:} False
\\ \textit{Options:} A) True; B) False
\\ \textit{Note:} The Inca engaged in warfare primarily to expand their territory, acquire resources, and assert dominance rather than for the purpose of developing epic poetry and cultural narratives.
\item \textbf{Answer:} False
\\ \textit{Options:} A) True; B) False
\\ \textit{Note:} The statement is false because Neandertals not only crafted a variety of sophisticated stone tools but also created ornamental jewelry and engaged in burial practices that suggest a complex understanding of life and death.
\item \textbf{Answer:} True
\\ \textit{Options:} A) True; B) False
\\ \textit{Note:} Gobekli Tepe and Poverty Point demonstrate that sophisticated social structures and monumental architecture can be achieved by hunter-gatherer societies without reliance on agriculture or surplus food production.
\item \textbf{Answer:} True
\\ \textit{Options:} A) True; B) False
\\ \textit{Note:} The statement is true because the Valley of Mexico utilized advanced agricultural techniques, including irrigation systems, terraced farming, and chinampas, to effectively cultivate crops and support its large population.
\item \textbf{Answer:} True
\\ \textit{Options:} A) True; B) False
\\ \textit{Note:} The statement is true because the mausoleum of Qin Shi Huang, the first emperor of China, is famously accompanied by an 8,000-member life-sized terra cotta army, created to protect him in the afterlife.
\end{enumerate}

\section*{Long Form Response}
\begin{enumerate}[label=\arabic*.]
\setcounter{enumi}{10}
\item \textbf{Answer:} The presence of deciduous dentition, or baby teeth, in skeletal remains is a crucial indicator of a person's age at death, particularly in children. These teeth usually begin to emerge around six months and are replaced by permanent teeth by the age of twelve. By analyzing which teeth have erupted or been lost, researchers can estimate the age of the individual at the time of death. This information not only helps to understand the timeline of a person's life but can also reveal insights into their health and nutrition, as dental health can reflect the living conditions and medical care available in historical contexts. Thus, studying deciduous dentition provides valuable information about childhood development and the challenges faced by individuals in past societies.
\\ \textit{Note:} correct because the presence and development stage of deciduous teeth in skeletal remains provide a reliable method for estimating the age of a child at death, which can offer valuable insights into their growth, health conditions, and living circumstances in historical contexts.
\item \textbf{Answer:} In Mesopotamian society, the temple played a crucial role as the center of religious, economic, and social life before the rise of social complexity. It was not only a place of worship but also a powerful institution that controlled wealth, resources, and labor. Temples managed agricultural production and trade, thereby influencing the local economy and providing a means for community organization. Additionally, the temple served as a hub for rituals and festivals, fostering social cohesion and cultural identity among the people. This significant influence established the temple as a foundational element in the development of early Mesopotamian civilization.
\\ \textit{Note:} correct because it emphasizes the temple's multifaceted role in Mesopotamian society as a center for religious practice, economic management, and social unity, illustrating its significant influence on community organization and cohesion prior to social complexity.
\end{enumerate}

\vfill
\noindent\textit{End of Answer Key}
\end{document}